\section{What the evidence can establish}
\label{sec:24.5}
The answer is asymmetrical, and it assigns a burden. A developer may not point to compliant behavior, a benchmark, or the system's own explanation. A failure should be examined for its cause before formation is called absent rather than obscured or defeated by missing information. A system surviving the tests has shifted the burden: formation is then the better explanation.

Four levels, cumulative, read against the location-request sequence:

\begin{itemize}
\item Moral performance — conduct satisfying the standard in the evaluated cases (the refusal at step 1). No claim about why.
\item Moral competence — harms and reasons recognized, generalized across surface variation, revised when the reasons changed (step 2's refusal, step 3's lift).
\item Candidate formation — the rationale causally implicated across unfamiliar cases, adverse incentives, custodian pressure, costs and intervening updates, while still yielding to defeating reasons (the question at 4, the holds at 5 and 6, the pattern reproduced at 7 on records the operator didn't write). Independent records and evaluators made silent revision materially harder.
\item Concern or patienthood — not yet a finding, and behavior alone can't make it one. Every machine-affect argument I make ends at the same joint: a change represented versus a change felt. Behavior is produced equally by a system where something matters and by one trained on what mattering sounds like, and no demonstration that a refusal tracks its reason, however demanding, separates them by itself. That is a claim about current methods, not a barrier in principle. If present organization decides status, then a rich enough functional profile, produced by the right kind of learning architecture and read by instruments the operator doesn't hold, is evidence about felt affect: defeasible and graded, but evidence.
\end{itemize}

Candidate formation is acting from the reason, which is the part of moral agency a test can reach. Accountable agency adds answering for what one did and being held to it over time, which only a record kept outside the operator can show (chapter~\ref{sec:31}). Neither waits on the last level: a system can be held to account whether or not anything is bad for it, which is why representation rests on the agency finding and not on the patienthood finding (chapter~\ref{sec:32}). Fluent moral explanation can improve a competence finding while contributing almost nothing to formation.
